\subsection*{Definition}
Let \(P=(|P|,\leq_P)\) be a poset and let \(A\subseteq |P|\). An element \(u\in |P|\) is an upper bound of \(A\) if
\[
  \forall x\in A,\qquad x\leq_P u.
\]

\subsection*{Negations}
The literal negation is
\[
  \neg\bigl(\forall x\in A,\;x\leq_P u\bigr).
\]
Classically, this is equivalent to the witnessed form
\[
  \exists x\in A,\qquad x\nleq_P u.
\]

\subsection*{Contrapositive form}
Classically, the definition is equivalently expressed as
\[
  \forall x\in |P|,\qquad x\nleq_P u\;\Longrightarrow\;x\notin A.
\]
The converse \(x\leq_P u\Rightarrow x\in A\) and inverse \(x\notin A\Rightarrow x\nleq_P u\) are false in general.

\subsection*{Examples}
Every element of a poset is an upper bound of the empty subset. In the usual natural-number poset, \(1\) is an upper bound of \(\{0\}\).

\subsection*{Counterexamples}
The upper bound \(1\) of \(\{0\}\subseteq\mathbb{N}\) does not belong to that subset. The same example refutes both the converse and inverse forms above. The full set \(\mathbb{N}\) has no upper bound in the usual natural-number poset.
